\section{Setting up and running existing models}

The most mature use of agents around existing models is to produce their input files, run them and repair the input when a run fails.

The closest analogue to preparing GGM's input data is AutoSAM, which builds input files for a nuclear reactor simulation code from engineering reports, diagrams and spreadsheets \citep{abulawi2026autosam}.
Two of its design choices carry over directly.
The agent first writes every extracted parameter to a human-readable intermediate file that an expert can check and correct before the model input is generated.
And where the documents are silent, it records the value it substituted as an assumption rather than filling the gap silently.
The extraction results show why the check is needed: 45 of 51 items expected from the PDF text were retrieved.

MOOSEnger, built around a multiphysics simulation code, shows how much depends on the system around the language model \citep{li2026moosenger}.
Its harness checks every candidate input against the code's schemas and runs the code's own validation before accepting it, and sends failures back to the model with the reason.
On 200 benchmark tasks, a frontier model alone produced an input that ran in 5\% of cases; inside the harness it did so in 89.5\%.
The authors stress that a simulation that runs is not necessarily right, since it can still encode the wrong equations or assumptions.

The same pattern appears in fluid dynamics, building energy and power systems \citep{yue2026foam, jiang2024eplus, jin2025gridmind}.
The language model writes the input, a solver or validator produces every number, and errors are fed back until the input passes, which establishes that the model runs but not that its result is correct.
